% generated by run/ -- do not edit by hand
\begin{tabular}{llrrrrrrllrrlr}
\toprule
 & & \multicolumn{3}{c}{Occurrence (\% points)} & \multicolumn{3}{c}{P(landscape $\mid$ recorded)} & \multicolumn{2}{c}{Affinity} & \multicolumn{2}{c}{Permutation importance} & \multicolumn{2}{c}{IndVal.g}\\
\cmidrule(lr){3-5}\cmidrule(lr){6-8}\cmidrule(lr){9-10}\cmidrule(lr){11-12}\cmidrule(lr){13-14}
Tribe & Species & Fpl & Pan & Cer & Fpl & Pan & Cer & BRF & RF$_{\mathrm{unw}}$ & Mean [95\% int.] & $p$ & Group & $p$\\
\midrule
Epiponini & \textit{Agelaia pallipes} (Olivier, 1791) & 72 & 0 & 0 & 0.69 & 0.15 & 0.16 & Fpl & Fpl & 0.092 [0.070, 0.123] & {0.050} & Fpl (0.85) & \textbf{$\leq$0.0001}\\
Epiponini & \textit{Polybia ruficeps} Schrottky, 1902 & 0 & 36 & 0 & 0.28 & 0.63 & 0.09 & Pan & Pan & 0.086 [0.044, 0.117] & \textbf{0.020} & Pan (0.60) & \textbf{0.002}\\
Epiponini & \textit{Agelaia bequaerti} (Richards, 1951) & 0 & 91 & 82 & 0.21 & 0.39 & 0.40 & Pan & Cer & 0.018 [-0.034, 0.063] & {0.400} & Pan (0.68) & \textbf{0.003}\\
Polistini & \textit{Polistes lanio} (Fabricius, 1775) & 33 & 0 & 0 & 0.53 & 0.22 & 0.25 & Fpl & Fpl & 0.013 [0.004, 0.025] & {0.300} & Fpl (0.58) & \textbf{0.006}\\
Epiponini & \textit{Polybia jurinei} de Saussure, 1854 & 50 & 0 & 0 & 0.59 & 0.20 & 0.21 & Fpl & Fpl & 0.011 [-0.003, 0.029] & {0.330} & Fpl (0.71) & \textbf{$<$0.001}\\
Epiponini & \textit{Polybia occidentalis} (Olivier, 1791) & 11 & 0 & 24 & 0.32 & 0.13 & 0.55 & Cer & Cer & 0.010 [-0.018, 0.044] & {0.350} & Cer (0.40) & {0.225}\\
Polistini & \textit{Polistes simillimus} Zikán, 1951 & 6 & 9 & 0 & 0.37 & 0.44 & 0.19 & Pan & Pan & 0.001 [-0.007, 0.008] & {0.230} & Pan (0.24) & {0.530}\\
Epiponini & \textit{Polybia sericea} (Olivier, 1792) & 22 & 0 & 0 & 0.43 & 0.26 & 0.30 & Fpl & Fpl & 0.000 [-0.001, 0.002] & {0.340} & Fpl (0.47) & {0.057}\\
Epiponini & \textit{Apoica flavissima} Van der Vecht, 1972 & 6 & 0 & 0 & 0.39 & 0.28 & 0.33 & rare & rare & 0.000 [0.000, 0.000] & {1.000} & Fpl (0.24) & {1.000}\\
Epiponini & \textit{Brachygastra lecheguana} (Latreille, 1824) & 6 & 0 & 0 & 0.39 & 0.28 & 0.33 & rare & rare & 0.000 [0.000, 0.000] & {1.000} & Fpl (0.24) & {1.000}\\
Mischocyttarini & \textit{Mischocyttarus consimilis} Zikán, 1949 & 0 & 0 & 6 & 0.34 & 0.25 & 0.41 & rare & rare & 0.000 [0.000, 0.000] & {1.000} & Cer (0.24) & {0.603}\\
Epiponini & \textit{Parachartergus smithii} (de Saussure, 1854) & 0 & 9 & 0 & 0.37 & 0.35 & 0.28 & rare & rare & 0.000 [0.000, 0.000] & {1.000} & Pan (0.30) & {0.241}\\
Polistini & \textit{Polistes versicolor} (Olivier, 1792) & 17 & 0 & 0 & 0.44 & 0.26 & 0.30 & Fpl & Fpl & -0.000 [-0.004, 0.006] & {0.640} & Fpl (0.41) & {0.103}\\
Epiponini & \textit{Polybia paulista} H. von Ihering, 1896 & 28 & 0 & 0 & 0.46 & 0.25 & 0.29 & Fpl & Fpl & -0.001 [-0.007, 0.002] & {0.480} & Fpl (0.53) & \textbf{0.023}\\
Epiponini & \textit{Apoica pallens} (Fabricius, 1804) & 22 & 0 & 6 & 0.48 & 0.19 & 0.34 & Fpl & Fpl & -0.004 [-0.015, 0.015] & {0.550} & Fpl (0.42) & {0.156}\\
Epiponini & \textit{Polybia chrysothorax} (Lichtenstein, 1796) & 0 & 9 & 6 & 0.37 & 0.34 & 0.30 & Pan & Pan & -0.032 [-0.054, -0.001] & {0.980} & Pan (0.23) & {0.706}\\
Epiponini & \textit{Polybia ignobilis} (Haliday, 1836) & 0 & 18 & 24 & 0.34 & 0.33 & 0.33 & Pan & Pan & -0.047 [-0.077, -0.018] & {0.960} & Cer (0.36) & {0.394}\\
\bottomrule
\end{tabular}
